\setlength{\parindent}{0pt}

\section{Teoria dos Jogos}

\textbf{Sprague-Grundy:} Todo jogo imparcial equivale a um Nim de um único monte. O valor de Grundy (nimber) de uma posição $u$ é:
\[g(u)=\text{mex}\{g(v): v \text{ alcançável de } u\}\]
onde $\text{mex}(S)$ é o menor inteiro não negativo fora de $S$.

Posição perdedora $\Leftrightarrow$ $g(u)=0$. Em jogo composto de partes independentes:
\[g = g_1 \oplus g_2 \oplus \cdots \oplus g_k\]

\textbf{Nim:} $k$ montes de tamanhos $a_1,\ldots,a_k$. Posição perdedora $\Leftrightarrow$ $a_1\oplus\cdots\oplus a_k=0$.

\textbf{Nim misère} (quem pegar o último \emph{perde}): subtraia 1 de cada pilha; se o XOR resultante $\neq 0$, a posição original é ganhadora.

\textbf{Wythoff:} 2 montes $(a,b)$, $a\leq b$. Movimentos: remover qualquer qtd de um monte, ou a mesma qtd dos dois. Posição perdedora $\Leftrightarrow$ $a=\lfloor\phi k\rfloor$, $b=\lfloor\phi^2 k\rfloor$ para algum $k\geq0$, onde $\phi=(1+\sqrt5)/2$. Equivalentemente: $k=a$, $b-a=\lfloor\phi k\rfloor$.
